\begin{titlepage}
  \begin{center}

  {\Huge AXIS\_STIMULATOR}

  \vspace{25mm}

  \includegraphics[width=0.90\textwidth,height=\textheight,keepaspectratio]{img/AFRL.png}

  \vspace{25mm}

  \today

  \vspace{15mm}

  {\Large Jay Convertino}

  \end{center}
\end{titlepage}

\tableofcontents

\newpage

\section{Usage}

\subsection{Introduction}

\par
This core contains two modules. A writer, and reader that should be placed on the
output, and input of the device under test. This will stream data through till
is has read all data. Then once all data has been written AND tlast is set to high
the writer module will end the simulation.


  dep:
    depend:
      - AFRL:utility:helper:1.0.0

##### Verilog include the files needed in the module you are using the function in (see src below for list of files and funtions).

`include "util_helper_math.vh"


\subsection{Dependencies}

\par
The following are the dependencies of the cores.

\begin{itemize}
  \item fusesoc 2.X
  \item iverilog (simulation)
  \item cocotb (simulation)
\end{itemize}

\subsubsection{fusesoc\_info Depenecies}
\begin{itemize}
\item dep
	\begin{itemize}
	\item AFRL:utility:helper:1.0.0
	\end{itemize}
\item dep\_tb
	\begin{itemize}
	\item AFRL:simulation:clock\_stimulator
	\item AFRL:utility:sim\_helper
	\end{itemize}
\item dep\_vpi
	\begin{itemize}
	\item AFRL:vpi:binary\_file\_io:1.0.0
	\end{itemize}
\end{itemize}


\section{Architecture}
\par
The project contains two modules master\_axis\_stimulus and slave\_axis\_stimulus. The master\_axis\_stimulus is used to take
input data from the slave axis interface (input) and write it to a file. Essentially it goes DUT\_MASTER to MASTER\_AXIS\_STIMULUS.
slave\_axis\_stimulus is used to read a file and push that data to the master axis interface (output). Essentially it goes
SLAVE\_AXIS\_STIMULUS to DUT\_SLAVE.
\par
This core uses a custom library for reading and writing files called vpi\_binary\_file\_io. This library provides multithreaded file reads
using a ring buffer between processes. The core will also puncture data according to its bit type. X/Z values are tossed if they are contained
in a byte.

\begin{itemize}
  \item \textbf{tm\_stim\_axis} Contains two modules master\_axis\_stimulus, slave\_axis\_stimulus.
\end{itemize}

Please see \ref{Module Documentation} for more information per target.

\section{Building}

\par
The all AXIS stimulator modules are written in Verilog 2001. They should synthesize in any modern FPGA software. The core comes as a fusesoc packaged core and can be
included in any other core. Be sure to make sure you have meet the dependencies listed in the previous section.

\subsection{fusesoc}
\par
Fusesoc is a system for building FPGA software without relying on the internal project management of the tool. Avoiding vendor lock in to Vivado or Quartus.
These cores, when included in a project, can be easily integrated and targets created based upon the end developer needs. The core by itself is not a part of
a system and should be integrated into a fusesoc based system. Simulations are setup to use fusesoc and are a part of its targets.

\subsection{Source Files}

\subsubsection{fusesoc\_info File List}
\begin{itemize}
\item src
	\begin{itemize}
	\item src/tm\_stim\_axis.v
	\end{itemize}
\item tb
	\begin{itemize}
	\item {'tb/tb\_axis.v': {'file\_type': 'verilogSource'}}
	\end{itemize}
\end{itemize}


\subsection{Targets} \label{targets}

\subsubsection{fusesoc\_info Targets}
\begin{itemize}
\item default
	\begin{itemize}
	\item[$\space$] Info: Default file set.
	\end{itemize}
\item sim
	\begin{itemize}
	\item[$\space$] Info: Default icarus sim.
	\end{itemize}
\item sim\_rand\_data
	\begin{itemize}
	\item[$\space$] Info: Use random data as a input.
	\end{itemize}
\item sim\_rand\_ready\_rand\_data
	\begin{itemize}
	\item[$\space$] Info: Use random data as a input, with random ready.
	\end{itemize}
\item sim\_8bit\_count\_data
	\begin{itemize}
	\item[$\space$] Info: Use counter data as a input.
	\end{itemize}
\item sim\_rand\_ready\_8bit\_count\_data
	\begin{itemize}
	\item[$\space$] Info: Use counter data as a input, with random ready.
	\end{itemize}
\end{itemize}


\subsection{Directory Guide}

\par
Below highlights important folders from the root of the directory.

\begin{enumerate}
  \item \textbf{docs} Contains all documentation related to this project.
    \begin{itemize}
      \item \textbf{manual} Contains user manual and github page that are generated from the latex sources.
    \end{itemize}
  \item \textbf{src} Contains source files for axis\_stimulator.
  \item \textbf{tb} Contains test bench files.
\end{enumerate}

\newpage

\section{Simulation}
\par
There is no simulation at the moment. Maybe a future addition?

\newpage

\section{Module Documentation} \label{Module Documentation}

\par
There project has multiple modules. The targets are the top system wrappers.

\begin{itemize}
\item \textbf{tm\_stim\_axis}
\item \textbf{tb\_axis}
\end{itemize}
The next sections document the module in great detail.

